\documentclass[11pt]{article}
\usepackage[spanish,mexico]{babel}
\usepackage[utf8]{inputenc}
\usepackage[T1]{fontenc}
\usepackage{geometry}
\usepackage{hyperref}
\usepackage{enumitem}
\geometry{margin=2.5cm}
\title{Estado del arte: motores epistemicos, RAG con grafos y fidelidad verificable}
\author{Sistema Operativo de Tesis de Posgrado}
\date{31 de mayo de 2026}
\begin{document}
\maketitle

\section*{Resumen}
El estado del arte apunta a pasar de recuperacion vectorial aislada a arquitecturas RAG gobernadas por conocimiento, grafos y evaluacion explicita. Para Toltecayotl, la consecuencia directa es que la fidelidad no debe depender de una sola respuesta generativa: requiere fuentes trazables, recuperacion estructurada, juez de calidad y bitacora auditable.

\section*{Hallazgos}
\begin{enumerate}[leftmargin=*]
\item Los surveys recientes de RAG reportan que la reduccion de alucinaciones depende tanto del recuperador como de metricas de evaluacion y trazabilidad del contexto.
\item GraphRAG y KG-RAG agregan relaciones factuales entre fragmentos; esto mejora coherencia y navegabilidad frente a ventanas de contexto planas.
\item La evaluacion epistemica madura combina fidelidad al contexto, consistencia logica, densidad de evidencia y registro reproducible de la decision.
\item El diseno operativo recomendado para la tesis es un motor hibrido: RAG local para privacidad, grafo de conocimiento para estructura y juez externo solo para contenido publico o redactado.
\end{enumerate}

\section*{Implicaciones para Toltecayotl}
El MCT debe operar como compuerta de calidad, no como validacion humana. Su salida debe incluir puntaje, hallazgos, modelo juez, identificador de solicitud y estado de revision humana. Esta separacion preserva el contrato de gobernanza: el sistema puede auditar, pero no validar autonomamente decisiones principales.

\section*{Referencias seleccionadas}
\begin{itemize}[leftmargin=*]
\item A Survey on Knowledge-Oriented Retrieval-Augmented Generation, arXiv:2503.10677.
\item Knowledge Graph-Guided Retrieval Augmented Generation, arXiv:2502.06864.
\item A Survey of Graph Retrieval-Augmented Generation for Customized Large Language Models, arXiv:2501.13958.
\item A Systematic Literature Review of Retrieval-Augmented Generation: Techniques, Metrics, and Challenges, MDPI Data 2025.
\end{itemize}
\end{document}
